% !PRINTED 33
\weberpage{33}
%<<CONT
drei Fällen gehen zwei ineinander über durch Vertauschung von
$V$, $\lambda^2$ mit $V:\lambda^2$, $1:\lambda^2$.

Wir wollen zwei für die Folge besonders wichtige unter diesen
Transformationen vollständig durchführen.

\begin{center}
1. \gesperrt{Die Gausssche Transformation.}
\end{center}

\[
V = \xi\eta, \qquad U = (a\xi+b\eta)^2.
\]

Wenn nun noch $U - \lambda^2 V$ ein Quadrat sein soll, so ist

\[
\lambda^2 = 4ab
\]

zu setzen, und es wird

\[
U - \lambda^2 V = (a\xi - b\eta)^2.
\]

Da $U-V$ sodann durch $\xi-\eta$ und durch $\xi-\varkappa^2\eta$ teilbar sein
muß, so ergeben sich, wenn man $\xi=\eta$, $\xi=\varkappa^2\eta$ setzt, aus
$U=V$ die Bedingungen

\[
a+b=\pm 1, \qquad a\varkappa^2+b=\pm\varkappa,
\]

woraus, wenn die oberen Zeichen genommen werden,

\[
a=\frac{1}{1+\varkappa}, \qquad b=\frac{\varkappa}{1+\varkappa}, \qquad
\lambda=\frac{2\sqrt{\varkappa}}{1+\varkappa},
\]

\[
V=\xi\eta, \qquad U=\left(\frac{\xi+\varkappa\eta}{1+\varkappa}\right)^2,
\qquad U-\lambda^2V=\left(\frac{\xi-\varkappa\eta}{1+\varkappa}\right)^2,
\]

\[
U-V=\frac{(\xi-\eta)(\xi-\varkappa^2\eta)}{(1+\varkappa)^2},
\]

\[
T=\frac{\xi^2-\varkappa^2\eta^2}{(1+\varkappa)^3}, \qquad
H=\frac{\xi^2-\varkappa^2\eta^2}{(1+\varkappa)^2},
\]

\[
M=\frac{1}{1+\varkappa}.
\]

Setzt man also wieder

\[
\frac{y}{x}=z, \qquad \frac{\eta}{\xi}=\zeta,
\]

so haben wir das Resultat, daß durch die Substitution

\weq{2}{z=\frac{(1+\varkappa)^2\zeta}{(1+\varkappa\zeta)^2}, \qquad
\lambda^2=\frac{4\varkappa}{(1+\varkappa)^2}}

die Transformation

\weq{3}{\frac{dz}{\sqrt{z(1-z)(1-\lambda^2z)}} =
\frac{(1+\varkappa)\,d\zeta}{\sqrt{\zeta(1-\zeta)(1-\varkappa^2\zeta)}}}

geleistet wird.
